% Research-note module, v21. Keep actionable material and exact scope only.
\subsection{\texorpdfstring{$G_{5,15}$}{G5,15} の継続に使う還元}
\label{ng51521:section}

\subsubsection{今回確定した境界}

低辺を $23,24$ とする標識では，core $2$ が二高，$3,4$ が三高，$0,1$ が四高である．
公開稿 \nolinkurl{g5_15_reduction_v21.tex} の定理 \texttt{g51521:reduction} は，
共通低 puncture $2$ を持つ共有高辺 $34$ の reverse 分岐を完全に分ける．
正例外型は無上界で排除された．通常型は中心二価 core の四点と十四の低 row 配置に限られる．
負例外型は $q_{34}=3,5$ に限られ，node 配分まで明示した二つの template が残る．
特に $q_{34}\ge7$ では，二価 core の二高分母が canonical 順で
$(7,31),(9,11),(31,7),(19,43)$ のいずれかになる．
支持 $G_{5,15}$ 自体は未排除であり，支持型の個数を減らしてはならない．

公開稿は全ての通常型中心，十四配置，正例外の八矛盾，負例外 template を保持する．
これらの表を研究ノートで重複させる必要はない．
検証器 \nolinkurl{verify_g5_15_reduction_v21.py} は旧版を import せず，
全 signed divisors と全 affine 半直線を処理する．
証明ファイルの SHA-256 と検証器自身の SHA-256 を通常版 JSON に記録する．

\subsubsection{通常型で次に解くべき五つの局所問題}

通常型では，共有高 ray への低 row 成分が正である三高 core を先に復元できる．
外部二標識の交換を除くと，低 row $R$ と低 $D$ の組は次の五つだけである．
\begin{center}\small
\begin{tabular}{c|c|r}
中心側の配置&$R$&低 $D$\\\hline
$M_1$, $|u|=4$&$(-1,-4,4)$&5\\
$M_1$, $u=1$&$(-1,-1,1)$&$-9$\\
$M_2$, $|u|=2$&$(-1,-2,2)$&2\\
$M_3$, $|u|=2$&$(-1,-2,2)$&$-7$\\
$M_4$, $|u|=2$&$(-1,-2,2)$&9
\end{tabular}
\end{center}
従って次の有望な目標は，これらの固定 $(R,D_0)$ に対する全正整数模型の解析である．
単に $q_{34}$ の探索上限を増やすより，この五問題の無上界分類または真に無限な族の構成を優先する．
もし三つ目と五つ目だけに有限個の模型が残ると証明できれば，
既存の固定 $q_{34}$ lifting 手順で通常型全体を解ける可能性がある．
有限範囲の疎な観測はその有限性の証明ではない．

計算補助として，$R=(-1,-m,m)$，低 column $(-S,-T,U)$ を取り，
$D_0=S+T-U-1=d$ を課す．高 rays の順を $a,b,c$ とし
\[
 q_a=S+1,\quad q_b=mT+1,\quad q_c=mU+1,
 \qquad U=S+T-d-1
\]
と置く．$H=p_{ab}>0$ なら $p_{ac}=-(H+1)$ である．
\[
 A=q_aq_b-1,\quad B=(m-1)T+1,\quad
 C=q_aq_c-1,\quad D=q_bq_c-1,\quad E=(m+1)U+1
\]
を用いると，残る高 row sum は
\begin{equation}\label{ng51521:simplified-quadratic}
 (D-EB)H^2+\{CB+D-E(A+B)\}H+A(C-E)=0
\end{equation}
に等しい．これは
\[
 \frac C{H+1}+\frac{DH}{A+BH}=E
\]
を払分母した式である．二次係数は
\[
 D-EB=(T-1)(U+1)
\]
と簡単になる．$T=1$ の境界は一次式に退化するため，この式だけで全域を genuine quadratic として扱ってはならない．
$H\mid A$，$H+1\mid C$，$A/H+B\mid D$，全符号，奇分母，および元の row/pair equations も必要である．
この単純化は今後の約数分解や整数不等式の出発点になる．

\subsubsection{固定共有分母の完全 lifting と有限探索の区別}

公開稿の命題 \texttt{g51521:fixed-Q-finite} によって，通常型の十四配置の各々と
任意の固定した奇数 $Q=q_{34}\ge3$ に対し，cores $2,3,4$ の全相互接着を有限に復元できる．
正 reverse 側で $m=|u|$，$U=(Q-1)/m$ とすると，低 $D$ が残る二 scale の和を固定する．
次に辺 $34$ の五つの node 分配を列挙し，反対側の共有高 row を決める．
反対側の最後の scale は共有高 $D$ の一次式で決まる．
一次式が共鳴しても，残る pair equation の二次係数
\[
 \frac{(m-1)^2U+m-2}{\{(m-1)U+1\}^2}>0
\]
が非零なので，高々二候補しか生じない．
この有限性は分母 $Q$ を固定した上での完全性であり，全 $Q$ に共通する上界ではない．

\nolinkurl{ordinary_fixed_q_lifting_v21.py} はこの手順を標準ライブラリのみで実装し，
各正局所模型に正の有理 weights の witness を付ける．例えば
\begin{verbatim}
python ordinary_fixed_q_lifting_v21.py --q 13 --output fixed_q13.json
\end{verbatim}
で，$q_{34}=13$ の十四配置を他の scale cutoff なしで全復元する．
出力に含まれるのは三 cores と三つの相互辺であり，四高 cores $0,1$ は含まれない．

診断用に
\begin{verbatim}
python ordinary_fixed_q_lifting_v21.py --scan-q 1001 --output q1001.json
\end{verbatim}
を実行した．$3\le Q\le1001$ の奇数五百個について，正 reverse 側の局所模型記録は十五件，
三 core 接着の完成は零件であった．十五件は十四配置への重複割当を含む．
この範囲で外部二標識の交換を除いた正 reverse 側模型は，
$R=(-1,-2,2)$ に対する次の四つである．scale 順は $(S,T,U)$ である．
\begin{center}\small
\begin{tabular}{r|rrr|r}
低 $D$&$S$&$T$&$U$&$Q$\\\hline
2&6&3&6&13\\
9&4&8&2&5\\
9&34&5&29&59\\
9&22&26&38&77
\end{tabular}
\end{center}
この表を全模型の分類として引用してはならない．
$Q>1001$ の未検証範囲を残した有界探索であり，公開稿の無上界排除は正例外型だけである．

\subsubsection{負例外で保持すべき小分母の構造}

負例外の共有高 ray は両側で低 ray と負に結ばれるため，
通常型のように正 reverse 側から二 scale の和を固定する手順をそのまま流用できない．
ただし $Q=3,5$ は強い制約である．公開稿の第二 template
$u=-(Q-1),v=-1$ では，共有高 rows は外部標識 $(0,1)$ の順で
\[
 (X,-(X+Q-1),-1),\qquad (-(Y+1),Y,-(Q-1))
\]
となる．辺 $34$ の puncture $2$ の node は二であるから，
$n_{34;0}+n_{34;1}=3$ であり，
\[
 X-Y=Q(1-n_{34;0})+1,\qquad 0\le n_{34;0}\le3
\]
を得る．第一 template $Q=3,u=-1,v=-2$ では rows は
\[
 (X,-(X+1),-2),\qquad (-(Y+2),Y,-1)
\]
であり，$X-Y=5-3n_{34;0}$ である．
これに低 $D$ と共有高 $D$ の相殺を加え，二つの局所 pair 系を同時に消去することが次の具体的な経路になる．
中心の正の枝は第二 template の $L=3$ だけに残り，
$T=S+2$，$r<S<2r$ を常に保持できる．
負の枝には $S=T=w$ と $w^2<2(r-w)(r-w+1)<4w^2$ を用いる．

\subsubsection{継続上の判断}

今回の進展を前提にするなら，通常型の五つの固定 $(R,D_0)$ 問題，
または負例外の $Q=3$ における共有高 row の連立消去が，次の作業を限定できる具体的な候補である．
$G_{5,14}$ の低辺補 graph $2K_2$ は今回の還元の対象外であり，未着手である．
支持全体に依存する $G_{5,15}$ の定理を，$G_{5,14}$ や $G_{5,16},G_{5,17}$ の部分 graph に適用してはならない．
新しい構造定理や構成例が別方向の価値を高める場合には，この優先順位を変更してよい．
